\documentclass[border=0pt]{standalone}
\usepackage[T1]{fontenc}
\usepackage[utf8]{inputenc}
\usepackage[spanish,es-noshorthands,es-nodecimaldot]{babel}
\usepackage{amsmath,amssymb}
\usepackage{newpxtext,newpxmath}
\usepackage{xcolor}
\usepackage{tikz}
\usetikzlibrary{arrows.meta,positioning,fit,backgrounds,calc}
\definecolor{navy}{HTML}{17243A}
\definecolor{blue}{HTML}{2A5CAA}
\definecolor{orange}{HTML}{D2691E}
\definecolor{cream}{HTML}{FCEFE3}
\definecolor{pale}{HTML}{EEF2F8}
\definecolor{text}{HTML}{1E232B}
\definecolor{muted}{HTML}{5F6873}
\definecolor{rule}{HTML}{C8D0DA}
\definecolor{sky}{HTML}{9DB9E3}
\definecolor{redact}{HTML}{3C4552}

\newcommand{\no}[1]{{\color{orange}#1}}
\setlength{\parindent}{0pt}
\setlength{\parskip}{0.45em}
\begin{document}
\begin{minipage}{26.783cm}\fontsize{20}{24.8}\selectfont\raggedright 
Distribución \textbf{observada} del desgaste al año siguiente en tres mercados.
Donde reparar es barato, más hogares reparan y la masa se corre a la izquierda. El
precio de reparación mueve cuántos reparan, no cómo reparar cambia el desgaste: es la
variable excluida que identifica la reparación (Hu \& Xin, 2024).
\end{minipage}
\end{document}
